% Dessin propre au CIR, séance 4 : le même circuit pile, interrupteur,
% lampe, interrupteur fermé puis ouvert.
% Les interrupteurs sont dessinés à la main : ceux de circuitikz (cspst,
% ospst) se distinguent à peine l'un de l'autre à cette taille.
% Auteur : Pierre F4JWI, 2026. Licence CC BY 4.0.
% Unités FIXÉES (x=1cm, y=1cm) : \DARCimage met un dessin à l'échelle en
% agrandissant les unités TikZ, pas les symboles de circuitikz ni le texte.
% Ici, le xshift (absolu) ne suivait pas : les deux circuits se
% chevauchaient (épreuve du 03/10/2026). Fixées, le dessin est composé à
% sa taille naturelle.
\begin{circuitikz}[thick, x=1cm, y=1cm]
  % interrupteur fermé : la lampe brille
  \draw (0,0) to[battery1, l=pile] (0,2.5) -- (0.85,2.5);
  \draw (0.85,2.5) -- (1.65,2.5) -- (2.5,2.5);
  \fill (0.85,2.5) circle (0.06) (1.65,2.5) circle (0.06);
  \node[above] at (1.25,2.6) {interrupteur};
  \draw (2.5,2.5) to[lamp, l_=lampe, name=L1] (2.5,0) -- (0,0);
  \foreach \a in {-60,-30,0,30,60} {
    \draw[DARCorange, very thick] ($(L1)+(\a:0.45)$) -- ($(L1)+(\a:0.75)$);
  }
  \node[align=center] at (1.25,-0.75) {interrupteur fermé\\la lampe brille};
  % interrupteur ouvert : la lampe est éteinte
  \begin{scope}[xshift=5.5cm]
    \draw (0,0) to[battery1, l=pile] (0,2.5) -- (0.85,2.5);
    \draw (0.85,2.5) -- ++(30:0.85);
    \draw (1.65,2.5) -- (2.5,2.5);
    \fill (0.85,2.5) circle (0.06) (1.65,2.5) circle (0.06);
    \node[above] at (1.25,2.95) {interrupteur};
    \draw (2.5,2.5) to[lamp, l_=lampe] (2.5,0) -- (0,0);
    \node[align=center] at (1.25,-0.75) {interrupteur ouvert\\la lampe est éteinte};
  \end{scope}
\end{circuitikz}
